\documentclass[Main.tex]{subfiles}
\begin{document}
\section{Đề kiểm tra định kì --- Mã đề 101}
\subsection{Bài tập}

%%%=================Phần trắc nghiệm nhiều lựa chọn===================%%%
\phan{Bài tập trắc nghiệm nhiều lựa chọn}
\textit{\large Thí sinh trả lời từ câu 1 đến câu 12. Mỗi câu thí sinh chỉ chọn một phương án}
\Opensolutionfile{ansex}[Ans/LGEX-C03_DE_KIEM_TRA_45P]
\Opensolutionfile{ans}[Ans/Ans-C03_DE_KIEM_TRA_45P]

%%%=============EX_1=============%%%
\begin{ex}%[6K3N1-1]
	Đồ vật sau đây làm từ vật liệu nào?
	\begin{center}
		\includegraphics[width=0.28\linewidth]{Images/DE_CUONG_HOA_6_HKI/dktra_binh_nuoc.png}
	\end{center}
	\choice
	{Kim loại}
	{Cao su}
	{\True Nhựa}
	{Thủy tinh}
	\loigiai{Bình đựng nước trong hình được làm từ nhựa: dẻo, nhẹ, không dẫn điện, dễ tạo hình.}
\end{ex}

%%%=============EX_2=============%%%
\begin{ex}%[6K3N1-1]
	Vật liệu có tính chất trong suốt là
	\choice
	{kim loại đồng}
	{\True thủy tinh}
	{gỗ}
	{thép}
	\loigiai{Thủy tinh trong suốt, cho ánh sáng đi qua nên được dùng làm cửa kính, chai lọ.}
\end{ex}

%%%=============EX_3=============%%%
\begin{ex}%[6K3H1-1]
	Đồ vật sau đây làm từ những vật liệu nào?
	\begin{center}
		\includegraphics[width=0.32\linewidth]{Images/DE_CUONG_HOA_6_HKI/dktra_dao.png}
	\end{center}
	\choice
	{\True Kim loại và gỗ}
	{Nhựa và kim loại}
	{Kim loại và thủy tinh}
	{Cao su và thủy tinh}
	\loigiai{Con dao gồm lưỡi dao làm bằng kim loại và cán dao làm bằng gỗ.}
\end{ex}

%%%=============EX_4=============%%%
\begin{ex}%[6K3N2-1]
	Vật thể nào sau đây được xem là nguyên liệu?
	\choice
	{Kim loại sắt (iron)}
	{Chậu nhựa}
	{\True Quặng bauxite}
	{Khí gas}
	\loigiai{Quặng bauxite là vật liệu tự nhiên chưa qua xử lí, cần chuyển hóa để sản xuất nhôm nên là nguyên liệu.}
\end{ex}

%%%=============EX_5=============%%%
\begin{ex}%[6K3H2-1]
	Khi dùng gỗ để sản xuất giấy thì người ta sẽ gọi gỗ là
	\choice
	{vật liệu}
	{\True nguyên liệu}
	{nhiên liệu}
	{phế liệu}
	\loigiai{Gỗ ở đây là vật liệu thô cần chuyển hóa để tạo ra sản phẩm là giấy nên được gọi là nguyên liệu.}
\end{ex}

%%%=============EX_6=============%%%
\begin{ex}%[6K3H3-3]
	Để sử dụng nhiên liệu tiết kiệm và hiệu quả cần phải cung cấp một lượng không khí hoặc oxygen
	\choice
	{\True vừa đủ}
	{thiếu}
	{dư}
	{tùy ý}
	\loigiai{Cung cấp lượng oxygen vừa đủ giúp nhiên liệu cháy hoàn toàn, tránh lãng phí và hạn chế sinh khí độc.}
\end{ex}

%%%=============EX_7=============%%%
\begin{ex}%[6K3N4-1]
	Trường hợp nào sau đây là thực phẩm?
	\choice
	{\True Quả nho}
	{Sắn}
	{Ngô}
	{Khoai}
	\loigiai{Quả nho là thực phẩm. Sắn, ngô, khoai chứa nhiều tinh bột nên là lương thực.}
\end{ex}

%%%=============EX_8=============%%%
\begin{ex}%[6K3H4-1]
	Trong các nhóm chất sau, những nhóm chất nào cung cấp năng lượng cho cơ thể?
	\begin{enumerate}[(1)]
		\item Chất đạm.
		\item Chất béo.
		\item Tinh bột, đường.
		\item Chất khoáng.
	\end{enumerate}
	\choice
	{$(1)$, $(2)$ và $(4)$}
	{$(2)$, $(3)$ và $(4)$}
	{$(1)$, $(2)$, $(3)$ và $(4)$}
	{\True $(1)$, $(2)$ và $(3)$}
	\loigiai{Chất đạm, chất béo và tinh bột, đường đều cung cấp năng lượng. Chất khoáng không cung cấp năng lượng mà giúp nâng cao hệ miễn dịch.}
\end{ex}

%%%=============EX_9=============%%%
\begin{ex}%[6K3N1-1]
	Các tính chất của kim loại bao gồm:
	\choice
	{\True Dẻo, dẫn điện, dẫn nhiệt, ánh kim, bị ăn mòn}
	{Trong suốt, dẫn nhiệt kém, cứng nhưng giòn, dễ vỡ}
	{Dẻo, không dẫn điện, dẫn nhiệt kém, không bị ăn mòn, dễ bị biến dạng nhiệt}
	{Đàn hồi, bền, không dẫn điện và nhiệt, không thấm nước, dễ cháy}
	\loigiai{Kim loại có ánh kim, dẻo, dẫn điện và dẫn nhiệt tốt, có thể bị ăn mòn (gỉ).}
\end{ex}

%%%=============EX_10=============%%%
\begin{ex}%[6K3N1-1]
	Ứng dụng của gốm, sứ là
	\choice
	{làm dây dẫn điện, nồi đun nấu, khung nhà cửa}
	{làm bình hoa, chai lọ, cửa kính}
	{\True làm chum vại, bát đĩa, chậu hoa}
	{làm lốp xe, gioăng cao su, đệm}
	\loigiai{Gốm, sứ không bị ăn mòn, dẫn nhiệt kém nên được dùng làm chum vại, bát đĩa, chậu hoa.}
\end{ex}

%%%=============EX_11=============%%%
\begin{ex}%[6K3H2-3]
	Khi khai thác quặng sắt, ý nào sau đây là không đúng?
	\choice
	{Khai thác tiết kiệm vì nguồn quặng có hạn}
	{Tránh làm ô nhiễm môi trường}
	{\True Nên sử dụng các phương pháp khai thác thủ công}
	{Chế biến quặng thành sản phẩm có giá trị để nâng cao hiệu quả kinh tế}
	\loigiai{Khai thác thủ công cho hiệu quả thấp, tốn công sức và không an toàn nên đây là ý không đúng.}
\end{ex}

%%%=============EX_12=============%%%
\begin{ex}%[6K3V3-3]
	Cho các việc làm sau:
	\begin{enumerate}
		\item[(a)] Sử dụng bếp gas phải có van an toàn, dùng xong phải khóa van lại.
		\item[(b)] Sử dụng bếp than trong phòng kín để đảm bảo lượng nhiệt không bị thất thoát.
		\item[(c)] Khi đi xe máy gặp đèn đỏ thời gian $100$ s thì không được tắt máy.
		\item[(d)] Tạo các lỗ trong viên than tổ ong để than cháy tốt hơn.
		\item[(e)] Vặn nhỏ lửa bếp gas khi thức ăn gần chín.
	\end{enumerate}
	Số việc làm nhằm sử dụng nhiên liệu an toàn, tiết kiệm là
	\choice
	{$2$}
	{\True $3$}
	{$4$}
	{$5$}
	\loigiai{Các việc làm đúng là (a), (d), (e). Việc (b) sai vì đốt than trong phòng kín sinh khí độc carbon monoxide. Việc (c) sai vì dừng đèn đỏ lâu nên tắt máy để tiết kiệm nhiên liệu. Vậy có $3$ việc làm.}
\end{ex}
\Closesolutionfile{ans}
\Closesolutionfile{ansex}
%\bangdapan{Ans-C03_DE_KIEM_TRA_45P}

%%%=================Phần trắc nghiệm đúng sai===================%%%
\phan{Bài tập trắc nghiệm đúng sai}
\textit{\large Thí sinh trả lời từ câu 1 đến câu 2. Trong mỗi ý a), b), c), d) ở mỗi câu thí sinh chọn đúng hoặc sai}
\Opensolutionfile{ansex}[Ans/LGTF-C03_DE_KIEM_TRA_45P]
\Opensolutionfile{ansbook}[Ansbook/AnsTF-C03_DE_KIEM_TRA_45P]
\Opensolutionfile{ans}[Ans/Tempt-C03_DE_KIEM_TRA_45P]
\setcounter{ex}{0}

%%%=============TF_1=============%%%
\begin{ex}%[6K3H1-1]
	Xét các phát biểu về nhựa và gốm, sứ.
	\choiceTF
	{\True Nhựa dẻo, nhẹ, không dẫn điện, dẫn nhiệt kém, không bị ăn mòn, dễ bị biến dạng nhiệt}
	{Nhựa thường được dùng làm bình hoa, chai lọ, cửa kính, …}
	{Gốm sứ không bị ăn mòn, dẫn nhiệt kém, hầu như không dẫn điện, cứng và khó vỡ}
	{\True Gốm sứ thường dùng làm chum vại, bát đĩa, chậu hoa, …}
	\loigiai{
		\begin{itemchoice}[T1,F2,F3,T4]
			\itemch Đây là các tính chất đặc trưng của nhựa.
			\itemch Sai vì nhựa ít khi làm bình hoa và không dùng làm cửa kính (cửa kính làm bằng thủy tinh).
			\itemch Sai vì gốm sứ giòn, dễ vỡ chứ không phải khó vỡ.
			\itemch Gốm sứ được dùng làm chum vại, bát đĩa, chậu hoa.
	\end{itemchoice}}
\end{ex}

%%%=============TF_2=============%%%
\begin{ex}%[6K3H3-6]
	Xét các phát biểu về an ninh năng lượng.
	\choiceTF
	{\True An ninh năng lượng là sự đảm bảo đầy đủ năng lượng dưới nhiều dạng khác nhau, ưu tiên các nguồn năng lượng sạch và giá thành rẻ}
	{Cần tăng cường sử dụng các nhiên liệu tái tạo như nhiên liệu sinh học, năng lượng mặt trời, dầu mỏ, …}
	{\True Hạn chế sử dụng các nguồn năng lượng không tái tạo như than đá, khí thiên nhiên}
	{\True Tăng cường sản xuất và sử dụng xe máy, ô tô chạy điện thay vì chạy xăng}
	\loigiai{
		\begin{itemchoice}[T1,F2,T3,T4]
			\itemch Đây chính là khái niệm an ninh năng lượng.
			\itemch Sai vì dầu mỏ là nguồn năng lượng không tái tạo.
			\itemch Than đá, khí thiên nhiên là nhiên liệu hóa thạch, cần hạn chế sử dụng.
			\itemch Xe chạy điện giúp giảm tiêu thụ xăng dầu và giảm ô nhiễm môi trường.
	\end{itemchoice}}
\end{ex}
\Closesolutionfile{ans}
\Closesolutionfile{ansbook}
\Closesolutionfile{ansex}
%\bangdapanTF{AnsTF-C03_DE_KIEM_TRA_45P}

%%%=================Phần trắc nghiệm trả lời ngắn===================%%%
\phan{Bài tập trắc nghiệm trả lời ngắn}
\textit{\large Thí sinh trả lời từ câu 1 đến câu 4}
\Opensolutionfile{ansex}[Ans/LGSA-C03_DE_KIEM_TRA_45P]
\Opensolutionfile{ansexh}[Ans/AnsSA-C03_DE_KIEM_TRA_45P]
\setcounter{ex}{0}

%%%=============SA_1=============%%%
\begin{ex}%[6K3H2-2]
	Thành phần chính của đá vôi là calcium carbonate. Phân tử khối của calcium carbonate bằng bao nhiêu?
	\shortans{$100$}
	\loigiai{
		Calcium carbonate có công thức CaCO$_3$.
		\\
		Phân tử khối: $40 + 12 + 16 \cdot 3 = 100$.
		\\
		\textbf{\textit{Đáp số:} $100$.}
	}
\end{ex}

%%%=============SA_2=============%%%
\begin{ex}%[6K3H1-1]
	Cho các vật dụng: $(1)$ bát ăn cơm, $(2)$ máy tính, $(3)$ bình hoa, $(4)$ cửa sổ, $(5)$ chậu hoa. Hãy sắp xếp các vật dụng thường làm bằng gốm sứ thành bộ số theo số thứ tự tăng dần (ví dụ $12$, $134$, …).
	\shortans{$135$}
	\loigiai{
		\begin{itemize}
			\item Vật dụng thường làm bằng gốm sứ gồm bát ăn cơm $(1)$, bình hoa $(3)$, chậu hoa $(5)$.
			\item Máy tính làm từ nhựa và kim loại; cửa sổ thường làm từ thủy tinh, gỗ hoặc kim loại.
		\end{itemize}
		\textbf{\textit{Đáp số:} $135$.}
	}
\end{ex}

%%%=============SA_3=============%%%
\begin{ex}%[6K3N3-1]
	Cho các nhiên liệu: $(1)$ than, $(2)$ gỗ, $(3)$ xăng, $(4)$ gas, $(5)$ cồn. Hãy sắp xếp các nhiên liệu ở trạng thái rắn thành bộ số theo số thứ tự tăng dần (ví dụ $12$, $134$, …).
	\shortans{$12$}
	\loigiai{
		\begin{itemize}
			\item Than $(1)$ và gỗ $(2)$ ở trạng thái rắn.
			\item Xăng, cồn ở trạng thái lỏng; gas ở trạng thái khí.
		\end{itemize}
		\textbf{\textit{Đáp số:} $12$.}
	}
\end{ex}

%%%=============SA_4=============%%%
\begin{ex}%[6K3H4-1]
	Cho các phát biểu sau:
	\begin{enumerate}[(1)]
		\item Tinh bột, đường là nguồn cung cấp năng lượng chính cho các hoạt động của cơ thể.
		\item Protein có nhiều trong rau xanh, củ quả tươi.
		\item Lipid là nguồn dự trữ năng lượng trong cơ thể và có tác dụng chống lạnh.
		\item Cơ thể cần một lượng lớn khoáng chất và vitamin để duy trì hoạt động.
		\item Thực phẩm dễ bị ôi thiu nên cần bảo quản cẩn thận như cho vào tủ lạnh, hút chân không, …
	\end{enumerate}
	Hãy sắp xếp các phát biểu đúng thành bộ số theo số thứ tự tăng dần (ví dụ $12$, $134$, …).
	\shortans{$135$}
	\loigiai{
		\begin{itemize}
			\item $(2)$ sai vì protein có nhiều trong cá, thịt, trứng, sữa và các loại đậu, đỗ.
			\item $(4)$ sai vì cơ thể chỉ cần một lượng nhỏ khoáng chất và vitamin.
			\item Các phát biểu đúng là $(1)$, $(3)$, $(5)$.
		\end{itemize}
		\textbf{\textit{Đáp số:} $135$.}
	}
\end{ex}
\Closesolutionfile{ansexh}
\Closesolutionfile{ansex}
%\bangdapanSA{AnsSA-C03_DE_KIEM_TRA_45P}
%%%=================Phần bài tập tự luận===================%%%
\phan{Bài tập tự luận}
\textit{\large Thí sinh trả lời từ câu 1 đến câu 3}
\Opensolutionfile{ansbth}[Ans/LGBT-C03_DE_KIEM_TRA_45P]
\Opensolutionfile{ansbt}[Ans/AnsBT-C03_DE_KIEM_TRA_45P]

%%%=============BT_1=============%%%
\begin{bt}%[6K3N4-1]
	\textit{($1$ điểm)} Cho các từ/cụm từ: \textit{lương thực, thực phẩm, bảo quản, tươi sống, chế biến}. Hãy chọn từ/cụm từ phù hợp điền vào chỗ trống để hoàn thành các phát biểu sau:
	\begin{enumerate}
		\item[(a)] Gạo, ngô, khoai, sắn là các loại ..(1).. chính ở Việt Nam.
		\item[(b)] Thịt, cá, tôm là các ..(2).. thường được dùng trong các bữa ăn hằng ngày. Chúng được ..(3).. để trở thành các món ăn.
		\item[(c)] Các thực phẩm ở dạng ..(4).. như thịt, cá cần được ..(5).. ở nhiệt độ thích hợp để bảo đảm an toàn cũng như tăng thời gian sử dụng.
		\item[(d)] Các ..(6).., ..(7).. cung cấp năng lượng và các chất thiết yếu giúp chúng ta có một cơ thể khỏe mạnh.
	\end{enumerate}
	\loigiai{
		\begin{itemize}
			\item (a) $(1)$ lương thực.
			\item (b) $(2)$ thực phẩm, $(3)$ chế biến.
			\item (c) $(4)$ tươi sống, $(5)$ bảo quản.
			\item (d) $(6)$ lương thực, $(7)$ thực phẩm.
		\end{itemize}
	}
\end{bt}

%%%=============BT_2=============%%%
\begin{bt}%[6K3V2-4]
	\textit{($1$ điểm)} Sơ đồ sau đây cho thấy cây mía có nhiều ứng dụng trong thực tế:
	\begin{center}
		\includegraphics[width=0.95\linewidth]{Images/DE_CUONG_HOA_6_HKI/so_do_cay_mia.png}
	\end{center}
	Trong sơ đồ trên, hãy cho biết đâu là vật liệu, nguyên liệu, nhiên liệu?
	\loigiai{
		\begin{itemize}
			\item \textbf{Vật liệu:} rỉ đường, đường ăn, bã mía, nước mía, giấy.
			\item \textbf{Nguyên liệu:} cây mía.
			\item \textbf{Nhiên liệu:} lá mía, rễ mía, bã mía.
		\end{itemize}
	}
\end{bt}

%%%=============BT_3=============%%%
\begin{bt}%[6K3C3-6]
	\textit{($1$ điểm)} Đọc đoạn thông tin sau và trả lời các câu hỏi.
	\begin{center}
		\textbf{MỘT SỐ LOẠI NHIÊN LIỆU CỦA TƯƠNG LAI}
	\end{center}
	Trong những năm tới, rất có thể bạn sẽ thường xuyên thấy những chiếc ô tô chạy bằng những loại nhiên liệu dưới đây.

	\textbf{\textit{Hydrogen.}} Các nhà sản xuất đang lên kế hoạch nạp hydrogen vào ô tô như các loại xăng dầu thông thường. Khi đó hydrogen sẽ chuyển hóa năng lượng hóa học thành điện và cung cấp cho hoạt động của chiếc xe. Tất cả những gì xe thải ra trong quá trình vận hành sẽ chỉ là nước.

	\textbf{\textit{Dầu diesel sinh học.}} Diesel sinh học là loại nhiên liệu được sản xuất từ dầu thực vật hay mỡ động vật để trở thành nhiên liệu cho xe. Nó được đánh giá là một nhiên liệu sạch với mức khí thải thấp hơn nhiều so với các loại nhiên liệu thông thường. Hơn nữa, vì được sản xuất từ các nguyên liệu rẻ, sẵn có như đậu tương nên diesel sinh học giúp các quốc gia giảm sự phụ thuộc vào nguồn dầu nhập khẩu.

	\textbf{\textit{Nhiên liệu pha ethanol.}} Thông thường, ethanol được sản xuất từ quá trình lên men của ngũ cốc như ngô. Đây là một nguồn nhiên liệu sạch và sản sinh khí nhà kính thấp hơn so với các loại khác. Ethanol được đưa vào xe sau khi đã pha trộn với xăng tùy theo từng nồng độ khác nhau. Nhiều quốc gia hiện nay đang sử dụng E85 với tỉ lệ pha trộn $85\%$ ethanol và $15\%$ xăng về thể tích.
	\begin{enumerate}
		\item[(a)] Vì sao hydrogen được coi là nhiên liệu không gây ô nhiễm môi trường?
		\item[(b)] Sử dụng các nhiên liệu như hydrogen, dầu diesel sinh học, … có lợi gì đối với an ninh năng lượng của mỗi quốc gia?
		\item[(c)] Xăng E90 có tỉ lệ $90\%$ ethanol và $10\%$ xăng về thể tích. Người ta phải thêm bao nhiêu lít ethanol vào $1$ lít xăng E85 để có xăng E90? (Giả sử không có hao hụt thể tích khi pha trộn.)
	\end{enumerate}
	\loigiai{
		\begin{itemize}
			\item (a) Xe chạy bằng nhiên liệu hydrogen chỉ thải ra nước, không gây ô nhiễm môi trường.
			\item (b) Các quốc gia sẽ có những nguồn năng lượng sạch, rẻ, bảo đảm nhu cầu sử dụng, giảm sự phụ thuộc vào dầu nhập khẩu.
			\item (c) Trong $1$ lít xăng E85 có $0{,}15$ lít xăng và $0{,}85$ lít ethanol.

			Thêm ethanol vào $1$ lít xăng E85 để có xăng E90 thì thể tích xăng không thay đổi và thể tích ethanol gấp $9$ lần thể tích xăng.

			Thể tích ethanol trong xăng E90 (sau khi được pha chế) là $0{,}15 \times 9 = 1{,}35$ (lít).

			Thể tích ethanol thêm vào là $1{,}35 - 0{,}85 = 0{,}5$ (lít).
		\end{itemize}
		\textbf{\textit{Đáp số:} thêm $0{,}5$ lít ethanol.}
	}
\end{bt}
\Closesolutionfile{ansbt}
\Closesolutionfile{ansbth}
\end{document}
